
\documentclass{article}
\usepackage{colortbl}
\usepackage{makecell}
\usepackage{multirow}
\usepackage{supertabular}

\begin{document}

\newcounter{utterance}

\centering \large Interaction Transcript for game `referencegame', experiment `line\_grids\_rows', episode 16 with reflexion{-}llama{-}3.1{-}8b{-}instruct.
\vspace{24pt}

{ \footnotesize  \setcounter{utterance}{1}
\setlength{\tabcolsep}{0pt}
\begin{supertabular}{c@{$\;$}|p{.15\linewidth}@{}p{.15\linewidth}p{.15\linewidth}p{.15\linewidth}p{.15\linewidth}p{.15\linewidth}}
    \# & $\;$A & \multicolumn{4}{c}{Game Master} & $\;\:$B\\
    \hline

    \theutterance \stepcounter{utterance}  
    & & \multicolumn{4}{p{0.6\linewidth}}{
        \cellcolor[rgb]{0.9,0.9,0.9}{
            \makecell[{{p{\linewidth}}}]{
                \texttt{\tiny{[P1$\langle$GM]}}
                \texttt{You have attempted this same task before and did not succeed. These are the plans you wrote after those attempts.} \\
\texttt{Trial \#1: New plan: I failed to distinguish the target grid from the distractors because my description focused on the rows of X's, which is common to all three grids. Next time, I will focus on the unique feature of the target grid, which is that it has a specific arrangement of X's and ▢'s. I will describe the target grid as the one that has a specific pattern of X's and ▢'s, such as "the grid with a ▢ in the top{-}left corner and a specific arrangement of X's in the first row".} \\
\\ 
\texttt{Trial \#2: New plan: I failed to distinguish the target grid from the distractors because my description focused on the rows of X's, which is common to all three grids. However, in my previous attempt, I also mentioned a specific arrangement of X's in the first row, which is not unique to the target grid. What I missed is that the target grid has a specific arrangement of X's and ▢'s, but I didn't specify what that arrangement is. Next time, I will focus on describing the unique pattern of X's and ▢'s in the target grid, such as "the grid with a ▢ in the top{-}left corner and a specific arrangement of X's in the first row, where the first row has X's in the second, third, and fourth positions".} \\
\\ 
\\ 
\texttt{You are given three grids, where each of them is 5 by 5 in size.} \\
\texttt{Grids have empty cells marked with "▢" and filled cells marked with "X".} \\
\texttt{Your task is to generate a referring expression that best describes the target grid while distinguishing it from the two other distractor grids.} \\
\texttt{The first grid is the target grid, and the following two grids are the distractors.} \\
\\ 
\texttt{Target grid:} \\
\\ 
\texttt{□ X X □ X} \\
\texttt{□ X X □ X} \\
\texttt{□ X X □ X} \\
\texttt{□ X X □ X} \\
\texttt{□ X X □ X} \\
\\ 
\texttt{Distractor grid 1:} \\
\\ 
\texttt{□ □ X X X} \\
\texttt{□ □ X X X} \\
\texttt{□ □ X X X} \\
\texttt{□ □ X X X} \\
\texttt{□ □ X X X} \\
\\ 
\texttt{Distractor grid 2:} \\
\\ 
\texttt{X X X X X} \\
\texttt{□ □ □ □ □} \\
\texttt{□ □ □ □ □} \\
\texttt{X X X X X} \\
\texttt{□ □ □ □ □} \\
\\ 
\texttt{Instruction: Describe the target grid.} \\
\texttt{Generate the referring expression starting with the tag "Expression: " for the given target grid. Omit any other text.} \\
            }
        }
    }
    & \\ \\

    \theutterance \stepcounter{utterance}  
    & \multicolumn{4}{p{0.6\linewidth}}{
        \cellcolor[rgb]{0.9,0.9,0.9}{
            \makecell[{{p{\linewidth}}}]{
                \texttt{\tiny{[P1$\rangle$GM]}}
                \texttt{Expression: the grid with a ▢ in the top{-}left corner and a specific arrangement of X's in the first row, where the first row has X's in the second, third, and fourth positions.} \\
            }
        }
    }
    & & \\ \\

    \theutterance \stepcounter{utterance}  
    & & & \multicolumn{2}{p{0.3\linewidth}}{
        \cellcolor[rgb]{0.9,0.9,0.9}{
            \makecell[{{p{\linewidth}}}]{
                \texttt{\tiny{[GM$|$GM]}}
                \texttt{Expression: the grid with a ▢ in the top{-}left corner and a specific arrangement of X's in the first row, where the first row has X's in the second, third, and fourth positions.} \\
            }
        }
    }
    & & \\ \\

    \theutterance \stepcounter{utterance}  
    & & \multicolumn{4}{p{0.6\linewidth}}{
        \cellcolor[rgb]{0.9,0.9,0.9}{
            \makecell[{{p{\linewidth}}}]{
                \texttt{\tiny{[GM$\rangle$P2]}}
                \texttt{You have attempted this same task before and did not succeed. These are the plans you wrote after those attempts.} \\
\texttt{Trial \#1: New plan: I selected the third grid because it seemed to have the most X's, but I didn't carefully read the expression. The expression specifically mentions four rows of five X's, with each row having one ▢. However, the third grid has only two rows of five X's, and the other rows have ▢'s. Next time, I will carefully read the expression and look for the exact pattern described. I will also consider the possibility that the grid with the most X's is not the correct answer.} \\
\\ 
\texttt{Trial \#2: New plan: I selected the first grid because it seemed to match the description of having a ▢ in the top{-}left corner, but I didn't pay attention to the specific arrangement of X's in the first row. The expression mentions that all rows have the same number of X's, but the first grid has a different number of X's in each row. Next time, I will carefully read the expression and look for the exact pattern described, including the number of X's in each row. I will also consider the possibility that the grid with the most X's is not the correct answer.} \\
\\ 
\\ 
\texttt{You are given three grids, where each of them is 5 by 5 in size.} \\
\texttt{Grids have empty cells marked with "▢" and filled cells marked with "X".} \\
\texttt{You are also given a referring expression that describes one of the given grids.} \\
\texttt{Your task is to select the grid that matches the given referring expression.} \\
\texttt{Generate only the number (in text) of the grid that matches the given expression by selecting first, second, or third.} \\
\\ 
\texttt{First:} \\
\\ 
\texttt{□ □ X X X} \\
\texttt{□ □ X X X} \\
\texttt{□ □ X X X} \\
\texttt{□ □ X X X} \\
\texttt{□ □ X X X} \\
\\ 
\texttt{Second:} \\
\\ 
\texttt{□ X X □ X} \\
\texttt{□ X X □ X} \\
\texttt{□ X X □ X} \\
\texttt{□ X X □ X} \\
\texttt{□ X X □ X} \\
\\ 
\texttt{Third:} \\
\\ 
\texttt{X X X X X} \\
\texttt{□ □ □ □ □} \\
\texttt{□ □ □ □ □} \\
\texttt{X X X X X} \\
\texttt{□ □ □ □ □} \\
\\ 
\texttt{Expression: the grid with a ▢ in the top{-}left corner and a specific arrangement of X's in the first row, where the first row has X's in the second, third, and fourth positions.} \\
\texttt{Question: Which grid does the expression refer to?} \\
\texttt{Start with the tag "Answer: ", followed by your selection. Omit any other text.} \\
            }
        }
    }
    & \\ \\

    \theutterance \stepcounter{utterance}  
    & & & \multicolumn{4}{p{0.6\linewidth}}{
        \cellcolor[rgb]{0.9,0.9,0.9}{
            \makecell[{{p{\linewidth}}}]{
                \texttt{\tiny{[GM$\langle$P2]}}
                \texttt{Answer: second} \\
            }
        }
    }
     \\ \\

    \theutterance \stepcounter{utterance}  
    & & & \multicolumn{2}{p{0.3\linewidth}}{
        \cellcolor[rgb]{0.9,0.9,0.9}{
            \makecell[{{p{\linewidth}}}]{
                \texttt{\tiny{[GM$|$GM]}}
                \texttt{second} \\
            }
        }
    }
    & & \\ \\

\end{supertabular}
}

\end{document}
